% Generated draft: requires mathematical and source review.
\documentclass[11pt]{article}
\usepackage[margin=25mm]{geometry}
\usepackage{fontspec}
\setmainfont{Latin Modern Roman}
\usepackage{amsmath,amssymb,mathtools,unicode-math}
\setmathfont{Latin Modern Math}
\usepackage{enumitem,needspace,xurl,hyperref}
\setlength{\parindent}{0pt}
\setlength{\parskip}{5pt}
\setlength{\emergencystretch}{3em}
\newcommand{\N}{\mathbb{N}}
\newcommand{\Z}{\mathbb{Z}}
\newcommand{\Q}{\mathbb{Q}}
\newcommand{\R}{\mathbb{R}}
\newcommand{\C}{\mathbb{C}}
\newcommand{\F}{\mathbb{F}}
\newcommand{\A}{\mathbb{A}}
\newcommand{\PP}{\mathbb{P}}
\newcommand{\E}{\mathbb{E}}
\newcommand{\eps}{\varepsilon}
\newcommand{\dd}{\,\mathrm{d}}
\begin{document}
\section*{UnsolvedMath ID: 30002338}
\textbf{Source problem code:} catalogue entry
\subsection{Definitions and mathematical statement}
For the Riemannian part of the question, write a sequence as \((M_j,g_j,\mu_j)_{j\in\N}\), where \((M_j,g_j)\) is a finite-dimensional Riemannian manifold and \(\mu_j\) is a measure to be specified. The condition that dimensions tend to infinity means
\[
 n_j:=\dim M_j,\qquad \lim_{j\to\infty}n_j=+\infty.
\]
The record does not prescribe whether \(\mu_j\) is Riemannian volume, normalized volume, or another measure. For the metric-measure part, write \((X_j,d_j,\mu_j)\), where \(d_j\) is a metric and \(\mu_j\) is a Borel measure. Neither the additional regularity of these spaces nor the notion of dimension used in the condition \(\dim X_j\to\infty\) is specified.

\textbf{Research question.} How should sequences of Riemannian or metric-measure spaces whose dimensions tend to infinity be studied, particularly with respect to concentration convergence, curvature-dimension conditions, and Laplacian eigenvalues?

The supplied question asks for a framework relating these three subjects, but does not specify a mathematical conclusion. In particular, it does not select a definition of concentration convergence, a limiting object, a curvature-dimension condition or its parameters, or a spectral convergence assertion. Consequently, it cannot yet be expressed as a uniquely determined quantified conjecture.

For clarity about the spectral terminology, if a nonnegative self-adjoint Laplacian realization \(L_j\) on \(L^2(X_j,\mu_j)\) has been specified, then an eigenvalue is a number \(\lambda\) for which
\[
 L_j u=\lambda u
 \quad\text{for some }u\in\operatorname{Dom}(L_j)\setminus\{0\}.
\]
This conditional definition does not assert that such a realization or any eigenvalues exist for the spaces in the record. Compactness, completeness, boundary conditions, discreteness of spectrum, metric and measure normalizations, and any scaling of eigenvalues remain unspecified. These choices are required before a precise convergence or stability statement can be formulated.
\subsection{Short English statement}
The question concerns spaces of increasing dimension and asks how concentration, curvature-dimension conditions, and Laplacian eigenvalues should be studied together. The supplied record presents a broad research program rather than a single precise conjecture; the class of spaces, convergence notions, curvature assumptions, and spectral target require specification.
\subsection{Sources}
\begin{itemize}
\item Local UnsolvedMath registry: \texttt{data/unsolvedmath/problems.json}, supplied record with permanent ID \(30002338\), title \emph{High-Dimensional Limits of Metric-Measure Spaces}, source problem code \texttt{catalogue entry}, in the Oberwolfach Reports Open Problems collection. No individual proposer is supplied; the proposed year is recorded as 2013.
\item Supplied report citation: \emph{Heat Kernels, Stochastic Processes and Functional Inequalities} (2014), Oberwolfach Reports, pp.~1359--1443, DOI \texttt{10.4171/owr/2013/23}. \url{https://doi.org/10.4171/owr/2013/23}. The background describes Oberwolfach Reports as documenting mathematical workshops and their open-problem discussions. No specific problem page or speaker attribution is supplied.
\item Supplied literature reference: Gigli--Mondino--Savare, \emph{Convergence of pointed non-compact metric measure spaces}. \url{https://doi.org/10.1112/plms/pdv047}.
\item Supplied literature reference: \emph{Oberwolfach Report 23/2013}. \url{https://doi.org/10.4171/owr/2013/23}.
\item Literature assessment dated \texttt{2026-08-21}: the supplied curation labels the problem open and states that convergence and stability theories already cover important cases of diverging dimension, while the broad infinite-dimensional concentration, Ricci, and spectral program is not a single settled theorem. This is supplied dated curation, not independently verified here; its time-sensitive status and cited claims require independent verification.
\item The supplied registry labels the statement \texttt{exact} and records a manual assessment that the extracted wording is self-contained, mathematically coherent, and free of LaTeX or report-layout artifacts. This is a supplied source-verification assessment, not an independent verification or a resolution of the missing formulation choices.
\end{itemize}
\end{document}
